\documentclass[12pt]{article}

\usepackage[margin=1in]{geometry}
\usepackage{setspace}
\usepackage{enumitem}
\usepackage{titlesec}

\setstretch{1.15}

\titleformat{\section}{\large\bfseries}{\thesection}{1em}{}
\titleformat{\subsection}{\normalsize\bfseries}{\thesubsection}{1em}{}
\titleformat{\subsubsection}{\normalsize\itshape}{\thesubsubsection}{1em}{}

\begin{document}
	
	This is the right moment to rethink the analytic structure, because the original analyses were built around a symmetric design that your data do not fully support, and the three-question reframing in the document you have shared is actually a better fit for both your theoretical framing and your results. Let me build a revised strategy that preserves that logic while being more specific about your data and more honest about where the predicted dissociation holds. \footnote{Source document: \texttt{Pasted markdown.md} :contentReference[oaicite:0]{index=0}}
	
	\section{Revised Analytic Strategy}
	
	\subsection{Orienting Principle}
	
	The three-question structure maps cleanly onto your introduction's theoretical argument. RTs establish \emph{who is faster} — they are a measure of efficiency, not morphological architecture. The N250 asks whether early form-based parsing is modulated by orthographic skill. The N400 asks whether later semantic integration is modulated by lexical-semantic knowledge. Individual differences are the diagnostic tool throughout, but their role shifts across questions: OS is foregrounded in Questions 1 and 2, LP is foregrounded in Question 3. This asymmetry should be explicit in how the models are structured and reported.
	
	\subsection{Handling Base Frequency and Family Size}
	
	This needs to be resolved before the models are specified, because the choice affects every analysis.
	
	\textbf{The problem:} Base Frequency and Family Size are moderately correlated at the item level, they appear together in multiple models, and your results suggest they do different things — Base Frequency drives RT effects for both words and nonwords, while Family Size drives ERP effects. Carrying both as independent predictors throughout every model produces multicollinearity and dilutes the theoretical signal.
	
	\textbf{Recommended solution: keep them separate but assign them to questions}
	
	Rather than collapsing them into a composite, assign each variable its primary analytic home based on what your data and theory suggest it indexes:
	
	\begin{itemize}[leftmargin=*]
		\item \textbf{Base Frequency} is your primary lexical familiarity index for \textbf{RT models} — it drives the behavioral effects and is interpretable as overall lexical availability. Family Size can be included as a covariate in RT models but should not be foregrounded.
		\item \textbf{Family Size} is your primary morphological structure index for \textbf{ERP models} — it drives the N250 and N400 effects and is interpretable as the richness of morphological family representation. Base Frequency can be included as a covariate in ERP models but should not be foregrounded.
	\end{itemize}
	
	This is not arbitrary — it is theoretically motivated by the distinction between token-level familiarity (Base Frequency) and type-level morphological structure (Family Size), and it is empirically supported by the pattern of your results. It also eliminates the high-order interactions that are currently making the models difficult to interpret.
	
	If a reviewer pushes back, the composite PCA approach is your fallback — but try the separation first, because it preserves the theoretical distinction that your introduction is built around.
	
	\subsection{Core Question 1: Does Orthographic Sensitivity influence overall lexical decision efficiency?}
	
	\textbf{Theoretical framing:} RTs index processing speed and overall efficiency. The finding that OS predicts overall RT but does not modulate the structure of morphological effects is theoretically meaningful — it suggests that orthographic skill affects how fast the system operates without changing what it does.
	
	\textbf{Model structure:}
	
	\textit{Words:}
	
	RT $\sim$ Base Frequency + Orthographic Sensitivity + (1 + Base Frequency $|$ participant) + (1 $|$ item)
	
	\begin{itemize}[leftmargin=*]
		\item Base Frequency as the primary lexical familiarity predictor
		\item Family Size included as a covariate but not interacted or foregrounded
		\item LP included once and reported as null — do not interact
	\end{itemize}
	
	\textit{Nonwords:}
	
	RT $\sim$ Complexity + Base Frequency + Orthographic Sensitivity + (1 + Complexity $|$ participant) + (1 $|$ item)
	
	\begin{itemize}[leftmargin=*]
		\item Complexity as the primary morphological contrast
		\item Base Frequency as lexical familiarity covariate
		\item LP included once and reported as null
	\end{itemize}
	
	\textbf{What to claim:} OS predicts overall efficiency across both lexical categories.
	
	\subsection{Core Question 2: Does Orthographic Sensitivity modulate early morpho-orthographic parsing (N250)?}
	
	\textbf{Theoretical framing:} The N250 indexes early form-based processing at the morpho-orthographic level.
	
	\textbf{Model structure:}
	
	\textit{Words (N250):}
	
	N250 $\sim$ Family Size $\times$ Orthographic Sensitivity + Language Proficiency + (1 $|$ participant:electrode) + (1 $|$ item)
	
	\textit{Nonwords (N250):}
	
	N250 $\sim$ Complexity $\times$ Orthographic Sensitivity + Complexity $\times$ Language Proficiency + Family Size + (1 $|$ participant:electrode) + (1 $|$ item)
	
	\textbf{What to claim:} Early morpho-orthographic parsing is modulated by orthographic skill.
	
	\subsection{Core Question 3: Does Language Proficiency shape semantic integration of morphological information (N400)?}
	
	\textbf{Theoretical framing:} The N400 indexes semantic integration.
	
	\textbf{Model structure:}
	
	\textit{Words (N400):}
	
	N400 $\sim$ Family Size $\times$ Language Proficiency + Orthographic Sensitivity + (1 $|$ participant:electrode) + (1 $|$ item)
	
	\textit{Nonwords (N400):}
	
	N400 $\sim$ Complexity $\times$ Family Size $\times$ Language Proficiency + Orthographic Sensitivity + (1 $|$ participant:electrode) + (1 $|$ item)
	
	\textbf{What to claim:} Language Proficiency shapes how morphological family structure is recruited during semantic integration.
	
	\subsection{Summary of Model Structure}
	
	\begin{center}
		\begin{tabular}{llll}
			\hline
			Question & Measure & Primary predictor & Primary moderator \\
			\hline
			Q1 Efficiency & RT Words & Base Frequency & OS \\
			Q1 Efficiency & RT Nonwords & Complexity + Base Frequency & OS \\
			Q2 Early parsing & N250 Words & Family Size & OS \\
			Q2 Early parsing & N250 Nonwords & Complexity & OS, LP \\
			Q3 Semantic integration & N400 Words & Family Size & LP \\
			Q3 Semantic integration & N400 Nonwords & Complexity $\times$ Family Size & LP \\
			\hline
		\end{tabular}
	\end{center}
	
	\subsection{What This Strategy Gains}
	
	The main advantage over the original analysis is that it replaces a symmetric design with an asymmetric design that reflects your theoretical predictions.
	
	\subsection{The Core Issue}
	
	Your rewritten introduction and the three-question analytic framework are well aligned, but the analyses were designed around an older symmetric structure.
	
	\subsection{What Needs to Change and Why}
	
	The models need to reflect asymmetry explicitly.
	
	\begin{itemize}[leftmargin=*]
		\item RT models: remove morphological $\times$ individual difference interactions
		\item N250 models: OS primary moderator
		\item N400 models: LP primary moderator
	\end{itemize}
	
	\subsection{What Does Not Need to Change}
	
	The introduction does not need substantial revision.
	
	\subsection{Recommended Order of Operations}
	
	\begin{enumerate}[leftmargin=*]
		\item Rerun RT models
		\item Rerun N250 models
		\item Rerun N400 models
		\item Rewrite results section
		\item Adjust aims if needed
	\end{enumerate}
	
\end{document}